\documentclass[12pt,a4paper]{article}
\usepackage[utf8]{inputenc}
\usepackage[T1]{fontenc}
\usepackage[spanish]{babel}
\usepackage{amsmath,amssymb}
\usepackage{geometry}
\usepackage{booktabs}
\usepackage{array}
\usepackage{xcolor}
\usepackage{tcolorbox}
\usepackage{fancyhdr}
\usepackage{lmodern}
\usepackage{microtype}
\usepackage{tikz}
\usetikzlibrary{arrows.meta, positioning, decorations.pathreplacing}

\geometry{top=2.5cm, bottom=2.5cm, left=2.5cm, right=2.5cm, headheight=25pt}

\definecolor{azulul}{RGB}{0,70,127}
\definecolor{verdeclaro}{RGB}{200,230,201}
\definecolor{azulclaro}{RGB}{200,220,240}
\definecolor{rojoclaro}{RGB}{255,210,210}
\definecolor{amarilloclaro}{RGB}{255,249,196}
\definecolor{grisclaro}{RGB}{245,245,245}

\tcbuselibrary{skins, breakable}

\pagestyle{fancy}
\fancyhf{}
\lhead{\color{azulul}\textbf{Econom\'ia Financiera 1 -- Gu\'ia 4}}
\rhead{\color{azulul}Ejercicio 4: Cia Minera del Zinc}
\cfoot{\thepage}
\renewcommand{\headrulewidth}{1pt}
\renewcommand{\headrule}{\hbox to\headwidth{\color{azulul}\leaders\hrule height \headrulewidth\hfill}}

% ---- Cajas de colores ----
\newtcolorbox{cajaenunciado}{
  colback=azulclaro!40, colframe=azulul, fonttitle=\bfseries,
  title=Enunciado del Ejercicio 4, breakable
}
\newtcolorbox{cajaintuicion}[1]{
  colback=verdeclaro!60, colframe=green!50!black, fonttitle=\bfseries,
  title=Intuici\'on: #1, breakable
}
\newtcolorbox{cajaresultado}{
  colback=verdeclaro!80, colframe=green!40!black, fonttitle=\bfseries\large,
  title=Resultado Final, halign title=center
}
\newtcolorbox{cajaformula}{
  colback=azulclaro!20, colframe=azulul!80, fonttitle=\bfseries,
  title=F\'ormula clave, breakable
}
\newtcolorbox{cajareflexion}{
  colback=amarilloclaro, colframe=orange!70!black, fonttitle=\bfseries,
  title=Reflexi\'on: \textquestiondown Por qu\'e el valor es casi exactamente US\$10?, breakable
}

\begin{document}

% ========================
%  PORTADA
% ========================
\begin{titlepage}
\centering
\vspace*{1cm}

{\color{azulul}\rule{\linewidth}{1.5pt}}
\vspace{0.5cm}

{\large\textbf{UNIVERSIDAD DE LIMA}\\[0.2cm]
Facultad de Ciencias Empresariales y Econ\'omicas\\
Carrera de Econom\'ia}

\vspace{1.5cm}

{\LARGE\color{azulul}\textbf{ECONOM\'IA FINANCIERA 1}}\\[0.4cm]
{\Large Gu\'ia de Pr\'actica 4 --- Mercado de Acciones}

\vspace{1.5cm}

{\color{azulul}\rule{0.6\linewidth}{0.8pt}}\\[0.5cm]
{\LARGE\textbf{Soluci\'on: Ejercicio 4}}\\[0.3cm]
{\Large\textit{Valuaci\'on de Cia Minera del Zinc}}\\[0.3cm]
{\Large\textit{Modelo de Descuento de Dividendos (DDM)}}\\[0.5cm]
{\color{azulul}\rule{0.6\linewidth}{0.8pt}}

\vspace{2cm}

\begin{tcolorbox}[colback=azulclaro!30, colframe=azulul, width=0.75\linewidth, halign=center]
\centering
\textbf{Datos clave del problema}\\[0.4cm]
\begin{tabular}{ll}
Dividendo base ($D_0$):   & US\$ 10{,}000{,}000 \\
Crecimiento a\~nos 1--3: & $g_1 = +8\%$         \\
Crecimiento a\~nos 4--6: & $g_2 = -10\%$        \\
Crecimiento a\~no 7+:    & $g_\infty = 0\%$     \\
Tasa de descuento:        & $r = 10\%$           \\
N\'umero de acciones:     & 10{,}000{,}000       \\
\end{tabular}
\end{tcolorbox}

\vfill
{\large Abril 2026}

{\color{azulul}\rule{\linewidth}{1.5pt}}
\end{titlepage}

% ========================
%  ENUNCIADO
% ========================
\begin{cajaenunciado}
Usted es analista de inversiones en un inversionista institucional y su jefe le pide que estime
el valor de \textbf{Cia Minera del Zinc}. Por estatutos sociales la empresa debe pagar
\textbf{100\% de sus utilidades en dividendos}. La empresa \emph{acaba de pagar} un dividendo
de \textbf{US\$10 millones}. Se espera que en los \textbf{pr\'oximos 3 a\~nos} el precio del
zinc aumente y se generen eficiencias operativas, aumentando la utilidad disponible para los
accionistas en \textbf{8\% cada a\~no}. A partir del \textbf{cuarto a\~no} la econom\'ia
mundial y especialmente la econom\'ia de China se desaceleraría y por lo tanto la utilidad
disponible para los accionistas \textbf{caer\'ia 10\% durante 3 a\~nos}, para luego
\textbf{estabilizarse (crecimiento cero) a perpetuidad}. Suponiendo que la mina tiene reservas
ilimitadas, \textquestiondown cu\'al es el \textbf{valor de la acci\'on} si el n\'umero de
acciones en circulaci\'on es \textbf{10 millones}? La tasa de descuento es \textbf{10\%}.
\end{cajaenunciado}

\vspace{0.5cm}

% ========================
%  INTUICIÓN GENERAL
% ========================
\begin{cajaintuicion}{El Modelo DDM y por qu\'e funciona aqu\'i}
El \textbf{Modelo de Descuento de Dividendos (DDM)} parte de un principio fundamental en
finanzas: \textit{el valor de un activo hoy es igual al valor presente de todos los flujos
de caja que ese activo generar\'a en el futuro}. Para una acci\'on, esos flujos son los
\textbf{dividendos}:

\[
P_0 = \sum_{t=1}^{\infty} \frac{D_t}{(1+r)^t}
\]

Como la empresa paga el \textbf{100\% de utilidades como dividendos}, los dividendos son
exactamente iguales a las utilidades por acci\'on. No hay retenci\'on, no hay reinversi\'on.
Esto simplifica el problema: valuar la empresa equivale a traer a valor presente todos sus
flujos de caja futuros.

El problema tiene \textbf{tres fases con tasas de crecimiento distintas}, por lo que no
podemos usar directamente la f\'ormula de Gordon (que supone $g$ constante para siempre).
En su lugar, calculamos los dividendos \textbf{a\~no por a\~no} durante el horizonte
expl\'icito (a\~nos 1--6), y usamos la f\'ormula de perpetuidad solo a partir del momento
en que el crecimiento se estabiliza (a\~no 7 en adelante).
\end{cajaintuicion}

\vspace{0.5cm}

\noindent Las tres fases son:

\begin{center}
\begin{tabular}{clcp{5.5cm}}
\toprule
\textbf{Fase} & \textbf{Per\'iodo} & \textbf{Tasa} & \textbf{Raz\'on econ\'omica} \\
\midrule
1 & A\~nos 1, 2, 3      & $g_1 = +8\%$       & Zinc sube de precio; eficiencias operativas \\
2 & A\~nos 4, 5, 6      & $g_2 = -10\%$      & Desaceleraci\'on de China \\
3 & A\~no 7 en adelante & $g_\infty = 0\%$   & Equilibrio de largo plazo (perpetuidad) \\
\bottomrule
\end{tabular}
\end{center}

\vspace{1cm}

% ========================
%  PASO 1
% ========================
\section{Paso 1: Calcular los dividendos a\~no por a\~no}

\begin{cajaintuicion}{Proyecci\'on de dividendos}
Sabemos que $D_0 = \$10\,\text{M}$ \emph{ya fue pagado}. Los dividendos futuros se obtienen
aplicando la tasa de crecimiento de cada fase:
\[
D_t = D_{t-1} \times (1 + g_{\text{fase}})
\]
Esta es la columna vertebral del DDM: construir el \emph{flujo de dividendos proyectados} antes
de descontarlos.
\end{cajaintuicion}

\subsection*{Fase 1: Crecimiento del 8\% --- A\~nos 1, 2 y 3}

La utilidad (y por tanto el dividendo) crece al 8\% por el auge del zinc:

\begin{align}
D_1 &= D_0 \times (1 + g_1) = 10{,}000{,}000 \times 1.08
    = \mathbf{\$10{,}800{,}000} \\[4pt]
D_2 &= D_1 \times (1 + g_1) = 10{,}800{,}000 \times 1.08
    = \mathbf{\$11{,}664{,}000} \\[4pt]
D_3 &= D_2 \times (1 + g_1) = 11{,}664{,}000 \times 1.08
    = \mathbf{\$12{,}597{,}120}
\end{align}

\noindent\textit{Comprobaci\'on directa:}
$D_3 = D_0 \times (1.08)^3 = 10{,}000{,}000 \times 1.259712 = \$12{,}597{,}120$. \checkmark

\subsection*{Fase 2: Ca\'ida del 10\% --- A\~nos 4, 5 y 6}

La desaceleraci\'on china reduce las utilidades un 10\% cada a\~no a partir de $D_3$:

\begin{align}
D_4 &= D_3 \times (1 + g_2) = 12{,}597{,}120 \times 0.90
    = \mathbf{\$11{,}337{,}408.00} \\[4pt]
D_5 &= D_4 \times 0.90 = 11{,}337{,}408.00 \times 0.90
    = \mathbf{\$10{,}203{,}667.20} \\[4pt]
D_6 &= D_5 \times 0.90 = 10{,}203{,}667.20 \times 0.90
    = \mathbf{\$9{,}183{,}300.48}
\end{align}

\subsection*{Tabla resumen de dividendos proyectados}

\begin{center}
\begin{tabular}{cccr}
\toprule
\textbf{A\~no} & \textbf{Fase} & \textbf{Tasa aplicada} & \textbf{Dividendo (US\$)} \\
\midrule
0 & Base (ya pagado) & ---    & 10{,}000{,}000.00 \\
1 & Expansi\'on      & $+8\%$ & 10{,}800{,}000.00 \\
2 & Expansi\'on      & $+8\%$ & 11{,}664{,}000.00 \\
3 & Expansi\'on      & $+8\%$ & 12{,}597{,}120.00 \\
4 & Contracci\'on    & $-10\%$& 11{,}337{,}408.00 \\
5 & Contracci\'on    & $-10\%$& 10{,}203{,}667.20 \\
6 & Contracci\'on    & $-10\%$&  9{,}183{,}300.48 \\
7+ & Perpetuidad     & $0\%$  &  9{,}183{,}300.48 (constante) \\
\bottomrule
\end{tabular}
\end{center}

\vspace{1cm}

% ========================
%  PASO 2
% ========================
\section{Paso 2: Calcular el Precio Terminal al final del A\~no 6}

\begin{cajaintuicion}{Capturar el valor de la perpetuidad}
A partir del a\~no 7, el dividendo es \textbf{constante para siempre} ($g_\infty = 0$). Sumar
infinitos t\'erminos uno por uno es imposible, pero existe una f\'ormula cerrada para una
\textbf{perpetuidad con crecimiento cero}: el valor al inicio de la perpetuidad es simplemente
el dividendo dividido entre la tasa de descuento.

En vez de descontar infinitos dividendos individualmente, ``colapsamos'' toda la perpetuidad en
un \textbf{precio terminal} $P_6$ que captura en un solo n\'umero el valor de todos los
dividendos desde el a\~no 7 en adelante, visto desde el a\~no 6. Luego ese $P_6$ se descuenta
como cualquier otro flujo.
\end{cajaintuicion}

\begin{cajaformula}
Si desde el per\'iodo $T+1$ los dividendos son constantes (g = 0) y la tasa de descuento es $r$:
\[
P_T = \frac{D_{T+1}}{r - g_\infty} = \frac{D_{T+1}}{r}
\]
En nuestro caso $T = 6$, y como $g_\infty = 0$: $D_7 = D_6 \times 1 = D_6$.
\end{cajaformula}

\begin{align}
D_7 &= D_6 \times (1 + 0) = 9{,}183{,}300.48 \\[8pt]
P_6 &= \frac{D_7}{r} = \frac{9{,}183{,}300.48}{0.10}
     = \mathbf{\$91{,}833{,}004.80}
\end{align}

\noindent\textbf{Interpretaci\'on:} $P_6 = \$91.83\,\text{M}$ representa el valor al final
del a\~no 6 de \emph{todos} los dividendos futuros desde el a\~no 7 en adelante (una
perpetuidad). Es como si al final del a\~no 6 pudi\'esemos vender la empresa a ese precio y
capturar todo el valor futuro.

\vspace{1cm}

% ========================
%  PASO 3
% ========================
\section{Paso 3: Calcular el Valor Presente de todos los flujos}

\begin{cajaintuicion}{Descontar es traer el futuro al presente}
El dinero en el futuro vale \textbf{menos} que el dinero hoy, porque si tuvi\'eramos ese dinero
hoy podr\'iamos invertirlo a la tasa $r = 10\%$ y obtener m\'as en el futuro. Descontar es
el proceso inverso: dado un flujo futuro, cuanto vale \emph{en este momento}.

El factor de descuento para el a\~no $t$ es $\frac{1}{(1+r)^t}$. Cuanto m\'as lejos est\'a
un flujo, menor es su contribuci\'on al valor presente, pues se divide entre un n\'umero m\'as
grande.
\end{cajaintuicion}

La f\'ormula completa del valor de la empresa es:

\[
V_0 = \underbrace{\sum_{t=1}^{6} \frac{D_t}{(1.10)^t}}_{\text{VP de dividendos expl\'icitos}}
    + \underbrace{\frac{P_6}{(1.10)^6}}_{\text{VP del valor terminal}}
\]

\subsection*{Factores de descuento $(1.10)^t$}

\begin{center}
\begin{tabular}{cc}
\toprule
\textbf{A\~no} $t$ & $(1.10)^t$ \\
\midrule
1 & 1.100000 \\
2 & 1.210000 \\
3 & 1.331000 \\
4 & 1.464100 \\
5 & 1.610510 \\
6 & 1.771561 \\
\bottomrule
\end{tabular}
\end{center}

\subsection*{C\'alculo detallado de cada Valor Presente}

\begin{align}
VP_1 &= \frac{10{,}800{,}000}{1.100000} = \mathbf{\$9{,}818{,}181.82} \\[6pt]
VP_2 &= \frac{11{,}664{,}000}{1.210000} = \mathbf{\$9{,}639{,}669.42} \\[6pt]
VP_3 &= \frac{12{,}597{,}120}{1.331000} = \mathbf{\$9{,}464{,}401.95} \\[6pt]
VP_4 &= \frac{11{,}337{,}408.00}{1.464100} = \mathbf{\$7{,}743{,}813.87} \\[6pt]
VP_5 &= \frac{10{,}203{,}667.20}{1.610510} = \mathbf{\$6{,}335{,}483.75} \\[6pt]
VP_6 &= \frac{9{,}183{,}300.48}{1.771561} = \mathbf{\$5{,}184{,}214.75} \\[10pt]
VP(P_6) &= \frac{91{,}833{,}004.80}{1.771561} = \mathbf{\$51{,}842{,}147.46}
\end{align}

\subsection*{Tabla resumen de valores presentes}

\begin{center}
\begin{tabular}{clrr}
\toprule
\textbf{A\~no} & \textbf{Flujo} & \textbf{Flujo (US\$)} & \textbf{Valor Presente (US\$)} \\
\midrule
1 & Dividendo $D_1$  & 10{,}800{,}000.00  &  9{,}818{,}181.82 \\
2 & Dividendo $D_2$  & 11{,}664{,}000.00  &  9{,}639{,}669.42 \\
3 & Dividendo $D_3$  & 12{,}597{,}120.00  &  9{,}464{,}401.95 \\
4 & Dividendo $D_4$  & 11{,}337{,}408.00  &  7{,}743{,}813.87 \\
5 & Dividendo $D_5$  & 10{,}203{,}667.20  &  6{,}335{,}483.75 \\
6 & Dividendo $D_6$  &  9{,}183{,}300.48  &  5{,}184{,}214.75 \\
6 & Precio terminal $P_6$ & 91{,}833{,}004.80 & 51{,}842{,}147.46 \\
\midrule
  & \textbf{TOTAL}  &                    & \textbf{100{,}027{,}912.02} \\
\bottomrule
\end{tabular}
\end{center}

\[
\boxed{V_{\text{empresa}} \approx \text{US\$}\;100{,}028{,}000}
\]

\vspace{1cm}

% ========================
%  PASO 4
% ========================
\section{Paso 4: Calcular el Valor por Acci\'on}

\begin{cajaintuicion}{De valor de empresa a valor por acci\'on}
El valor total de la empresa calculado ($\approx\$100\,\text{M}$) es el valor de
\emph{todos} los flujos de caja para \emph{todos} los accionistas. Para obtener el precio
por acci\'on, simplemente dividimos entre el n\'umero total de acciones en circulaci\'on.
\end{cajaintuicion}

\[
P_{\text{acci\'on}} = \frac{V_{\text{empresa}}}{N_{\text{acciones}}}
= \frac{100{,}027{,}912}{10{,}000{,}000}
\]

\begin{tcolorbox}[colback=verdeclaro!80, colframe=green!40!black,
                  halign=center, boxsep=8pt]
\Large
\[
\boxed{P_{\text{acci\'on}} \approx \textbf{US\$ 10.00 por acci\'on}}
\]
\end{tcolorbox}

\vspace{1cm}

% ========================
%  RESULTADO FINAL
% ========================
\begin{cajaresultado}
\begin{center}
\large
\begin{tabular}{p{7cm}p{5.5cm}}
\textbf{Valor total de la empresa:} & US\$ 100{,}027{,}912 $\approx$ US\$ 100 millones \\[4pt]
\textbf{Acciones en circulaci\'on:} & 10{,}000{,}000 acciones \\[4pt]
\textbf{Valor por acci\'on:}        & \textbf{US\$ 10.00 por acci\'on} \\
\end{tabular}
\end{center}
\end{cajaresultado}

\vspace{1cm}

% ========================
%  REFLEXIÓN
% ========================
\begin{cajareflexion}
Observemos qu\'e ocurrir\'ia si la empresa \emph{no} tuviese las fases de expansi\'on y
contracci\'on, sino que desde el inicio pagase dividendos constantes con $g = 0$:

\[
P_0^{*} = \frac{D_1}{r} = \frac{D_0 \times (1+0)}{r}
= \frac{\$10{,}000{,}000}{0.10} = \$100{,}000{,}000
\quad\Rightarrow\quad \$10 \text{ por acci\'on}
\]

El resultado es \emph{el mismo} que el que obtuvimos con las tres fases. \textquestiondown Por
qu\'e? Porque las fases de crecimiento ($+8\%$ durante 3 a\~nos) y de contracci\'on
($-10\%$ durante 3 a\~nos) se \textbf{compensan casi perfectamente en valor presente}.

Matem\'aticamente, el dividendo estabilizado al final del a\~no 6 es:
\begin{align*}
D_6 &= D_0 \times (1.08)^3 \times (0.90)^3
     = 10{,}000{,}000 \times 1.259712 \times 0.729 \\
    &= 10{,}000{,}000 \times 0.91833
     \approx 9{,}183{,}300
\end{align*}
Que es \emph{ligeramente inferior} a $D_0$, de ah\'i que el valor final sea
\$100.03M y no exactamente \$100M. La diferencia es m\'inima
($\approx 0.03\%$) porque $(1.08)^3 \approx 1.2597$ y $(0.90)^3 \approx 0.729$,
cuyo producto es $0.9183 \approx 1$ pero no exactamente 1.

\textbf{Lecci\'on clave del DDM:} lo que determina el valor de largo plazo de una empresa
es el \emph{nivel estabilizado} de sus dividendos y la tasa de descuento. Las fluctuaciones
transitorias (auges y recesiones) tienen un impacto secundario, especialmente cuando son
relativamente sim\'etricas en magnitud y duraci\'on.
\end{cajareflexion}

\vspace{1cm}

% ========================
%  LÍNEA DE TIEMPO
% ========================
\newpage
\section*{L\'inea de Tiempo del Problema}

\vspace{0.5cm}

\begin{center}
\resizebox{\linewidth}{!}{%
\begin{tikzpicture}[scale=0.95, font=\small]

% ---- Eje temporal ----
\draw[thick, -{Stealth}] (-0.3, 0) -- (14.5, 0) node[right] {Tiempo};

% ---- Marcas en el eje ----
\foreach \x/\t in {0/0, 1.8/1, 3.6/2, 5.4/3, 7.2/4, 9.0/5, 10.8/6, 12.6/7+} {
  \draw[thick] (\x, 0.12) -- (\x, -0.12);
  \node[below=4pt] at (\x, 0) {$t=\t$};
}

% ---- Fases (rectángulos sombreados) ----
\fill[green!25] (0, 0.5) rectangle (5.4, 1.6);
\draw[green!60!black, thick] (0, 0.5) rectangle (5.4, 1.6);
\node[align=center] at (2.7, 1.05) {\textbf{Fase 1}\\$g_1=+8\%$};

\fill[red!20] (5.4, 0.5) rectangle (10.8, 1.6);
\draw[red!60!black, thick] (5.4, 0.5) rectangle (10.8, 1.6);
\node[align=center] at (8.1, 1.05) {\textbf{Fase 2}\\$g_2=-10\%$};

\fill[blue!15] (10.8, 0.5) rectangle (14.2, 1.6);
\draw[blue!60!black, thick] (10.8, 0.5) rectangle (14.2, 1.6);
\node[align=center] at (12.5, 1.05) {\textbf{Fase 3}\\$g_\infty=0\%$};

% ---- Flechas de dividendos hacia abajo ----
\foreach \x/\d in {1.8/10.8M, 3.6/11.66M, 5.4/12.60M, 7.2/11.34M, 9.0/10.20M, 10.8/9.18M} {
  \draw[blue!70!black, -{Stealth}, thick] (\x, 0) -- (\x, -1.2);
  \node[below=2pt, blue!70!black, align=center] at (\x, -1.3) {\$\d};
}

% ---- D0 ----
\draw[gray, -{Stealth}] (0, 0) -- (0, -1.0);
\node[below=2pt, gray, align=center] at (0, -1.1) {\$10M\\(pagado)};

% ---- P6 ----
\draw[purple!80!black, -{Stealth}, very thick] (10.8, 0) -- (10.8, -2.3);
\node[below=2pt, purple!80!black, align=center] at (10.8, -2.4)
  {$P_6=\$91.83M$\\(perpetuidad)};

% ---- Etiqueta a la derecha ----
\node[right=4pt, blue!70!black] at (14.2, -0.7) {$D_{7+}=\$9.18M$};
\node[right=4pt, blue!70!black] at (14.2, -1.1) {(constante)};

% ---- Corchete "Descuentan a hoy" ----
\draw[thick, decorate, decoration={brace, amplitude=8pt, mirror}]
  (-0.1, -2.9) -- (10.9, -2.9);
\node[below=10pt, align=center] at (5.4, -2.9) {Todos estos flujos se descuentan\\a $t=0$ usando el factor $(1.10)^t$};

\end{tikzpicture}}
\end{center}

\vspace{0.8cm}

% ========================
%  VERIFICACIÓN NUMÉRICA
% ========================
\section*{Verificaci\'on num\'erica compacta}

\begin{align*}
V_0 &= \frac{10.8}{1.1^1} + \frac{11.664}{1.1^2} + \frac{12.597}{1.1^3}
     + \frac{11.337}{1.1^4} + \frac{10.204}{1.1^5}
     + \frac{9.183 + 91.833}{1.1^6}
     \quad \text{(millones US\$)} \\[6pt]
    &= 9.818 + 9.640 + 9.464 + 7.744 + 6.335 + 5.184 + 51.842
     = \mathbf{100.027 \text{ millones}}
\end{align*}

\[
\Rightarrow\quad P_{\text{acci\'on}} = \frac{100.027\text{ M}}{10\text{ M acciones}}
\approx \boxed{\textbf{US\$ 10.00}}
\]

\end{document}
